% !TEX root = ../../../main.tex
% Block card: catalogue.shape_cluster -- "Shape clustering" (Part I, the blocks chapter).
\cardfield{What it does}
Builds a family tree of the pool's training windows by shape, cuts it into a number of
clusters, and lets you name each cluster and mark it interesting or not.

\cardfield{When to use it}
In the research question 1 chain: \blockref{preprocessing.window_pool} $\to$
\blockref{preprocessing.trace_shape} $\to$ this block, then a model trained on its clusters.
It is the ``let the computer group the windows first'' arm, against the human-label arm of
\blockref{catalogue.manual_labels}.

\begin{plainterms}[A tree, and where to cut it]
Ward clustering starts with every window on its own and repeatedly joins the two groups whose
merger adds the least spread, until one group is left. The record of those joins is a tree
(a \emph{dendrogram}). Cutting the tree at some height gives clusters: cut low for many small
clusters, high for a few big ones. No rule says where the true cut is; the block proposes one
(the cut whose clusters are best separated, by a score called the silhouette, ignoring tiny
``speck'' clusters) and you can move it.
\end{plainterms}

\cardfield{What goes in, what comes out}
In go the shapes from Trace shape. Out comes a cluster number for every training window; every
validation, test and exam window gets $-1$, because only a model trained on the clusters may
assign them. The tree is kept on disk, so moving the cut later does not rebuild it.

\cardfield{Settings that matter}
\emph{k} is the number of clusters at the cut; 0 takes the proposal. \emph{sample} is how many
training windows build the tree (20,000); the others join the nearest cluster.
\emph{mapping}, set in the table on the block's page, names each cluster and marks it
interesting or not.

\cardfield{Watch out for}
Once a model trained on this pool has a test score, the cut and the mapping are frozen: a
different cut is refused. The proposal is a guide, not an answer: the choice of cut is yours.
